%%
%% This is file `sample-sigconf.tex',
%% generated with the docstrip utility.
%%
%% The original source files were:
%%
%% samples.dtx  (with options: `all,proceedings,bibtex,sigconf')
%% 
%% IMPORTANT NOTICE:
%% 
%% For the copyright see the source file.
%% 
%% Any modified versions of this file must be renamed
%% with new filenames distinct from sample-sigconf.tex.
%% 
%% For distribution of the original source see the terms
%% for copying and modification in the file samples.dtx.
%% 
%% This generated file may be distributed as long as the
%% original source files, as listed above, are part of the
%% same distribution. (The sources need not necessarily be
%% in the same archive or directory.)
%%
%%
%% Commands for TeXCount
%TC:macro \cite [option:text,text]
%TC:macro \citep [option:text,text]
%TC:macro \citet [option:text,text]
%TC:envir table 0 1
%TC:envir table* 0 1
%TC:envir tabular [ignore] word
%TC:envir displaymath 0 word
%TC:envir math 0 word
%TC:envir comment 0 0
%%
%% The first command in your LaTeX source must be the \documentclass
%% command.
%%
%% For submission and review of your manuscript please change the
%% command to \documentclass[manuscript, screen, review]{acmart}.
%%
%% When submitting camera ready or to TAPS, please change the command
%% to \documentclass[sigconf,anonymous]{acmart} or whichever template is required
%% for your publication.
%%
%%
\documentclass[sigconf]{acmart}

% SVG figures (for \includesvg)
\usepackage{svg}
\usepackage{float} % for [H] placement


%%
%% \BibTeX command to typeset BibTeX logo in the docs
\AtBeginDocument{%
  \providecommand\BibTeX{{%
    Bib\TeX}}}

%% Rights management information.  This information is sent to you
%% when you complete the rights form.  These commands have SAMPLE
%% values in them; it is your responsibility as an author to replace
%% the commands and values with those provided to you when you
%% complete the rights form.
%\setcopyright{acmlicensed}
%\copyrightyear{2026}
%\acmYear{2026}
%\acmDOI{XXXXXXX.XXXXXXX}
%% These commands are for a PROCEEDINGS abstract or paper.
%%%%%내꺼\acmConference[ASIA CCS '26]{ACM Asia Conference on Computer and Communications Security}{June 01--05, 2026}{Bangalore, India}
%%
%%  Uncomment \acmBooktitle if the title of the proceedings is different
%%  from ``Proceedings of ...''!
%%
%%\acmBooktitle{Woodstock '18: ACM Symposium on Neural Gaze Detection,
%%  June 03--05, 2018, Woodstock, NY}
%\acmISBN{978-1-4503-XXXX-X/2018/06}
\setcopyright{cc}
\setcctype[4.0]{by-nc-nd}
\acmYear{2026}\copyrightyear{2026}
\acmConference[ASIA CCS '26]{ACM Asia Conference on Computer and Communications Security}{June 1--5, 2026}{Bangalore, India}
\acmBooktitle{ACM Asia Conference on Computer and Communications Security (ASIA CCS '26), June 1--5, 2026, Bangalore, India}
\acmDOI{10.1145/3779208.3804888}
\acmISBN{979-8-4007-2356-8/26/06}

%%
%% Submission ID.
%% Use this when submitting an article to a sponsored event. You'll
%% receive a unique submission ID from the organizers
%% of the event, and this ID should be used as the parameter to this command.
%%\acmSubmissionID{123-A56-BU3}

%%
%% For managing citations, it is recommended to use bibliography
%% files in BibTeX format.
%%
%% You can then either use BibTeX with the ACM-Reference-Format style,
%% or BibLaTeX with the acmnumeric or acmauthoryear sytles, that include
%% support for advanced citation of software artefact from the
%% biblatex-software package, also separately available on CTAN.
%%
%% Look at the sample-*-biblatex.tex files for templates showcasing
%% the biblatex styles.
%%

%%
%% The majority of ACM publications use numbered citations and
%% references.  The command \citestyle{authoryear} switches to the
%% "author year" style.
%%
%% If you are preparing content for an event
%% sponsored by ACM SIGGRAPH, you must use the "author year" style of
%% citations and references.
%% Uncommenting
%% the next command will enable that style.
%%\citestyle{acmauthoryear}


%%
%% end of the preamble, start of the body of the document source.
\begin{document}

%%
%% The "title" command has an optional parameter,
%% allowing the author to define a "short title" to be used in page headers.
\title{POSTER: GEM-PLI Side Channel in GPON: Inferring Encrypted Streaming Traffic}

% Anonymous submission: keep author block commented for camera-ready.
%
%
\author{Hyunjun An}
\affiliation{%
  \institution{Heungdeok High School}
  \city{Yongin}
  \country{Republic of Korea}
}
\email{hjahn@anhyunjun.com}
\orcid{0009-0002-3984-4385}
\renewcommand{\shortauthors}{Hyunjun An}


\begin{abstract}
Gigabit Passive Optical Networks (GPON), an ITU-T standardized technology adopted in FTTH networks, supports downstream payload encryption to protect subscribers' traffic. Unfortunately, encryption does not cover all access layer metadata. This poster presents a preliminary analysis demonstrating that the Payload Length Indicator (PLI), which remains in plaintext in GEM headers, can be exploited as a side channel to infer streaming activity. Using a GPON downstream simulator, we demonstrate that PLI traces are highly consistent across repeated runs of the same Youtube video and suggest that a passive attacker on the same optical splitter as the victim could potentially infer a victim's streaming activity (e.g., Youtube video identification).
\end{abstract}

%%HJ\ccsdesc{Do Not Use This Code~Generate the Correct Terms for Your Paper}

\begin{CCSXML}
<ccs2012>
   <concept>
       <concept_id>10003033.10003039</concept_id>
       <concept_desc>Networks~Network protocols</concept_desc>
       <concept_significance>500</concept_significance>
       </concept>
   <concept>
       <concept_id>10003033.10003039.10003044</concept_id>
       <concept_desc>Networks~Link-layer protocols</concept_desc>
       <concept_significance>500</concept_significance>
       </concept>
 </ccs2012>
\end{CCSXML}

\ccsdesc[500]{Networks~Network protocols}
\ccsdesc[500]{Networks~Link-layer protocols}

\keywords{GPON, encrypted traffic analysis, side-channel inference, video identification, passive optical networks, protocol security}

%% (Removed sample teaser/received blocks; content-only change.)

\maketitle

\section{Introduction And Motivation}
As fiber internet has become dominant in network infrastructures, GPON--one of the passive optical network (PON) technologies--is adopted by many ISPs to provide fiber-to-the-home (FTTH) internet. Since GPON is a popular PON technology, it is important to ensure traffic confidentiality. 
 GPON supports AES-based downstream payload encryption at the GPON Encapsulation Method (GEM) layer. However, some fields in the GEM header remain plaintext such as the Payload Length Indicator (PLI) field used for multiplexing and the Port-ID field used for mapping subscribers' Optical Network Terminals/Units (ONTs/ONUs) (typically called "modem").

Such plaintext fields may appear innocuous. However, recent studies have demonstrated that such metadata may be exploited to infer a victim's activity (e.g., identifying a victim's web-based video streaming watching activity) ~\cite{bae-sec22}~\cite{ytyt}.

While prior studies primarily focused on upper layer or cellular traffic analysis~\cite{bae-sec22}~\cite{ytyt} , we researched whether the plaintext PLI field in GEM headers can be exploited to infer a victim's streaming activity. 
Also a limitation in many prior optical-communication security studies is that an adversary must physically install an inline optical tap at some point along the link to observe the victim's traffic. In contrast, GPON downstream traffic is broadcast over a passive splitter; thus, an attacker who is connected to the same splitter can receive downstream transmissions without deploying an additional inline tap.

In this poster, we developed a software-based GPON downstream simulator to analyze the statistical properties of PLI traces. Using these PLI traces, we present preliminary evidence that Youtube streaming traces are highly consistent over time and distinctive at the GEM layer, suggesting that the value of the PLI field can be exploited to infer a victim's video streaming activity.
\section{Background}  

\noindent\textbf{GPON framing.}
In GPON, the Optical Line Terminal (OLT) broadcasts downstream traffic over a passive optical splitter to multiple Optical Network Units (ONUs) (typically called "modem"). Downstream transmission is organized into fixed duration GTC frames (125~$\mu$s), each composed of a physical control block downstream (PCBd) followed by a downstream GTC payload region.
~\cite{itu-g9843} % Fig. 8-1 / 8-2

\noindent\textbf{GEM and plaintext metadata.}
Upper layer traffic (e.g., Ethernet frame) is encapsulated in the Payload field in GEM frames. A GEM frame starts with a 5-byte GEM header containing a 12-bit PLI, a 12-bit Port-ID, a 3-bit Payload Type Indicator (PTI), and a header error control (HEC) field. PTI indicates whether the payload is the end of a fragmented payload or not. %
~\cite{itu-g9843} % Fig. 8-11 + PLI/PTI text

\noindent\textbf{PLI length even under encryption.}
In GPON, downstream user traffic is carried as GEM frames multiplexed in the GTC payload~\cite{itu-g9843}.
Each GEM header includes a plaintext PLI field that specifies the payload length $L$ in bytes (up to 4095 bytes). In certain cases, larger payload values are fragmented accordingly~\cite{itu-g9843} (where PTI is used to indicate whether a GEM payload is an end of a fragment or not).
Downstream GEM encryption only encrypts payload field (not the GEM header) using AES-CTR.
Since AES-CTR encryption is length-preserving, even when payload bytes are encrypted, an attacker can observe the plaintext's payload length.

\noindent\textbf{Video Identification.} HTTP Adaptive Streaming (HAS) is the popular video streaming model -- which is generally adopted by video streaming platforms (e.g., Youtube) --delivers video by letting the client periodically request short media segments at a bitrate chosen by an adaptive bitrate (ABR) algorithm. This request--response process forms characteristic bursty download patterns and bitrate-dependent segment sizes, which can be exploited for traffic inference.~\cite{bae-sec22}

\section{Threat Model}
We consider an attacker who shares the same optical splitter as the victim.
\begin{itemize}
  \item \textbf{1. Premise} The attacker connects an optical capture device (e.g., an FPGA with GPON optics
or a modified ONU capable of logging raw GTC frames) to the splitter and
passively records the downstream broadcast from the OLT. The attacker does not need to register their device with the OLT because an attacker only needs downstream GTC frames; thus, there is simply no need to transmit any upstream traffic.
We assume the attacker can perform the necessary physical-layer processing to recover GTC frames (clock recovery, frame delineation, de-scrambling, and FEC decode) with sufficient lock time and bit error rate for header parsing. The attacker may employ commercial GPON capture equipment
(e.g., \textit{PON Analyzer by MT2 Company~\cite{MT2}} or \textit{GPONDOCTOR by EXXN Engineering~\cite{EXXN}}) --- devices originally designed for troubleshooting GPON infrastructures --- which claim to support raw GTC frame acquisition.
  \item \textbf{2. Collect GEMs} Then, attacker extracts GTC frames from downstream traffic and parses GEM frames. 
  \item \textbf{3. Convert GEM PLI to dataset} Then, the attacker converts parsed GEM headers into time-stamped PLI sequences.
  Since GEM headers carry \texttt{port\_id}, the attacker can organize the dataset \emph{per Port-ID} (i.e., build one PLI time series for each observed \texttt{port\_id} on the splitter).

  \emph{Optional:} if the attacker additionally knows (or infers via a side channel) the Port-ID used by a specific victim's ONU, they can filter the dataset to that \texttt{port\_id} and perform victim-specific inference.
  \item \textbf{4. Dataset format} Each record includes \texttt{t\_start\_us}, \texttt{gem\_id}, \texttt{port\_id}, and \texttt{pli}.
  \item \textbf{5. Classification via database matching} The attacker matches the captured PLI time series against a pre-collected database of PLI traces derived from known videos, selecting the closest match to identify the victim's video.
\end{itemize}
\begin{figure}[!htbp]
  \centering
  \includesvg[
    width=\linewidth,
    inkscapearea=drawing,
    inkscapelatex=false
  ]{fig/pip_with_screenshot_compact_readable_v5}
  \caption{Attack pipeline for PLI-based video identification.}
  \label{fig:pipeline}
  \Description{Attack pipeline for PLI-based video identification.}
\end{figure}


\section{Trace Generation \& Dataset}
To demonstrate the threat model in our experimental environment, we developed a software-based GPON downstream trace generator~\cite{github} that follows key ITU-T G.984.3 TC-layer parameters relevant to PLI (125~$\mu$s framing and GEM header formatting). From a downlink ethernet packet capture file of encrypted video traffic, we (i) map Ethernet frames into GEM payloads (fragmenting to satisfy the 12 bit PLI limit), (ii) generate plaintext GEM headers (PLI/Port-ID/PTI), and (iii) pack the GEM stream into fixed-length 125~$\mu$s downstream frames consistent with the configured line rate. The attacker signal is either the GEM-level PLI sequence or the per-frame occupancy series. Large Ethernet packets may generate multiple GEM fragments; thus, PLI values in our traces reflect per fragment payload lengths (the attacker-visible observable).

%\textbf{Note.} Payload bytes can optionally be AES-CTR protected; however, this does not change the PLI trace because AES-CTR is length-preserving (XOR), so encryption does not change any fragment length. Our analysis uses only plaintext headers (e.g., PLI) and timing.
\section{Experiment and Preliminary Results}
We played 10 YouTube videos (three runs each) on Ubuntu 24.04 at a fixed quality (720p) and captured 180 seconds of downlink traffic per run (PCAP) with a $\pm 3\,s$ buffer. We then replayed each PCAP in our simulator to generate a time-stamped PLI trace (one trace per run) and visualized the cumulative PLI--time curves, which are highly consistent across repeated runs of the same video (Figure 3).

For similarity analysis, we bin PLI events into 100ms bins, compute the cumulative sum, normalize by the final cumulative value, and compare traces using Pearson correlation and RMSE.
These traces approximate the payload length observations available to a real same-splitter adversary and enable systematic evaluation of PLI based inference.
We use run1/run2 dataset files as a database and run3 dataset files as query traces.

\begin{figure}[!t]
  \centering
  \includegraphics[width=0.85\linewidth]{fig/box_rmse.png}
  \includegraphics[width=0.85\linewidth]{fig/box_pearson_r.png}
  \caption{Two stacked box plots: RMSE (lower is better), Pearson r (higher is better).}
  \label{fig:pli-metrics}
  \Description{RMSE and Pearson Graph of PLI-Time}
\end{figure}

\begin{table}[!b]
  \centering
  \small
  \setlength{\tabcolsep}{3pt}
  \renewcommand{\arraystretch}{1.0}
  \caption{Per-query Top-1 predictions~\cite{github}.}
  \Description{A table listing, for each query video trace (run3), the top-1 predicted matching reference trace (run1/run2), whether it is correct, and the similarity scores ($r$/RMSE).}
  \label{tab:top1-detail}
  \begin{tabular}{l l c c c}
    \hline
    \textbf{Query} & \textbf{Prediction} & \textbf{CORRECT} & \textbf{margin} & \textbf{$r$/RMSE} \\
    \hline
    vid001/r3 & vid001/r1 & T & 2 & 0.998/0.025 \\
    vid002/r3 & vid002/r1 & T & 5 & 0.995/0.041 \\
    vid003/r3 & vid003/r2 & T & 1 & 0.980/0.050 \\
    vid004/r3 & vid004/r1 & T & 2 & 0.996/0.037 \\
    vid005/r3 & vid005/r2 & T & 2 & 0.998/0.020 \\
    vid006/r3 & vid008/r1 & F & 2 & 0.994/0.048 \\
    vid007/r3 & vid007/r1 & T & 2 & 0.999/0.025 \\
    vid008/r3 & vid008/r1 & T & 2 & 0.999/0.018 \\
    vid009/r3 & vid009/r2 & T & 2 & 0.990/0.041 \\
    vid010/r3 & vid010/r1 & T & 3 & 0.998/0.020 \\
    \hline
  \end{tabular}
  % --- insert right after \end{tabular} ---

  \vspace{0.15em}
  {\footnotesize\noindent\textbf{Note.} Top-1 is chosen by rank-voting over Pearson ($r\uparrow$) and RMSE ($\downarrow$).

  Margin is the best--runner-up gap (vote: rank-sum gap $\Delta$).
  A rare error likely stems from transient jitter/buffering that makes two PLI slopes similar.}\par
\end{table}

\begin{figure}[!t]
  \centering
  \includegraphics[width=0.9\linewidth]{fig/cumsum_vid001.pdf}
  \caption{Cumulative PLI - Time graph for one of 10 videos}
  \label{fig:plitime}
  \Description{Cummulative PLI vs TIME Graph}
\end{figure}

We quantify similarity on the normalized cumulative PLI time series: same video pairs exhibit higher Pearson correlation (median $r=0.991$ vs $0.970$, AUC=0.779) and lower RMSE (median 0.048 vs 0.098, AUC=0.826).
Using nearest matching with rank voting over Pearson and RMSE, we achieve 90\% Top-1 identification accuracy (9/10) when matching run3 queries against run1, run2 query traces (Table~\ref{tab:top1-detail}).

\section{Conclusion And Discussion}
Despite GEM Payload Encryption, GPON may remain vulnerable to streaming inference because GEM-header PLI values can be exploited to infer a victim’s video activity.

\textbf{Mitigations.} Potential mitigations include (i) padding / length-obfuscation within GPON framing and (ii) traffic shaping / distribution smoothing to reduce discriminative PLI patterns, at the cost of bandwidth/latency overhead. 

\textbf{Limitations:} 
\begin{itemize}
  \item \textbf{Simulator-driven:} vendor-specific DBA/scheduling and fragmentation/padding may change PLI dynamics.
  \item \textbf{No physical validation yet:} results are derived from a software simulator fed by PCAP traces; we have not yet validated end-to-end leakage on a commercial OLT/ONU/splitter.
  \item \textbf{Capture feasibility assumptions:} we assume the attacker can reliably recover downstream GTC/GEM headers via receive-only capture and physical-layer processing (e.g., clock recovery and FEC decode).
  \item \textbf{Robustness:} short windows, start-time shift, ABR, and cross-traffic are not exhaustively tested.
  \item \textbf{Scale/open-set:} larger catalogs and unknown queries may reduce accuracy; rejection thresholds are future work.
\end{itemize}

\textbf{Characteristics Unique to GPON:} 


\textbf{Future Work.}
We will validate PLI-based video-streaming identification on a physical GPON testbed with a commercial OLT/ONU. Using an FPGA-based capture device, we will extract downstream GTC/GEM headers and derive PLI time series.
We will also employ a CNN-based classifier and extend the evaluation to larger video catalogs, longer-term measurements, and adaptive streaming dynamics (e.g., resolution and bitrate changes). We hypothesize that XGS-PON has the same issue; the XGEM frame used in XGS-PON has a Payload Length Indicator field similar to the PLI field in GEM, and we plan to validate this on an XGS-PON testbed.

\bibliographystyle{ACM-Reference-Format}
\bibliography{references}

\end{document}
\endinput
%%
%% End of file `sample-sigconf.tex'.